\documentclass[11pt,a4paper]{article}
\usepackage[margin=23mm]{geometry}
\usepackage[T1]{fontenc}
\usepackage[utf8]{inputenc}
\usepackage{lmodern,microtype,amsmath,amssymb,booktabs,tabularx,array}
\usepackage{graphicx,xcolor,listings,enumitem}
\usepackage[section]{placeins}
\usepackage[colorlinks=true,linkcolor=blue!50!black,urlcolor=blue!50!black]{hyperref}
\usepackage{fancyhdr}
\pagestyle{fancy}
\fancyhf{}
\lhead{Statistical learning: worked notebook explanations}
\rfoot{\thepage}
\setlength{\headheight}{14pt}
\setlength{\parindent}{0pt}
\setlength{\parskip}{6pt}
\setlength{\emergencystretch}{2em}
\setlist{itemsep=3pt,topsep=4pt}
\lstset{language=Python,basicstyle=\small\ttfamily,breaklines=true,
 columns=fullflexible,keepspaces=true,showstringspaces=false,
 frame=single,rulecolor=\color{gray!35},backgroundcolor=\color{gray!4},
 commentstyle=\color{green!35!black},keywordstyle=\color{blue!65!black},
 aboveskip=8pt,belowskip=8pt}
\newcommand{\code}[1]{\texttt{\detokenize{#1}}}
\newcommand{\RSS}{\operatorname{RSS}}
\newcommand{\RMSE}{\operatorname{RMSE}}
\newcommand{\MAE}{\operatorname{MAE}}
\newcommand{\E}{\mathbb{E}}
\newcommand{\ind}[1]{\mathbf{1}\{#1\}}
\newcommand{\fig}[2]{\begin{figure}[htbp]\centering
\includegraphics[width=\linewidth,height=0.39\textheight,keepaspectratio]{figures/#1}
\caption{#2}\end{figure}}
\newcommand{\notebooktitle}[3]{\begin{center}{\LARGE\bfseries #1\par}\vspace{5pt}
{\large Problem statement and worked resolution\par}\vspace{5pt}
Companion to \code{#2}\par\end{center}
\textbf{How to use this explanation.} #3 Section numbers follow the notebook's numbered sections where present. Numerical values are rounded and refer to the recorded run, not universal performance claims.}
\newcommand{\methods}[1]{\subsection*{Main Python calls}#1}

\begin{document}
\notebooktitle{4C: USArrests, scaling, and descriptive groups}{04_C_usarrests.ipynb}{This explanation treats preprocessing as a choice of geometry, distinguishes direction coefficients from correlations, and qualifies what a historical full-data analysis can establish.}
\section*{Problem statement and scope}
Describe 50 historical US state records using Murder, Assault, UrbanPop, and Rape.
The three arrest-rate variables are measured per 100,000 residents for 1973; UrbanPop is a historical urban-population percentage.
The state names identify rows and are excluded from numerical features.
We compare principal component analysis (PCA) in original and standardized units, reconstruct a two-component approximation,
and form descriptive groups. All 50 rows participate; there is no target and no held-out prediction claim.

The notebook uses the conventional textbook table unchanged. The
\href{https://stat.ethz.ch/R-manual/R-devel/library/datasets/html/USArrests.html}{R dataset documentation}
notes a transcription error in Maryland's UrbanPop value and rounding issues in some other states.
In particular, the conventional Maryland entry is 67, while the documented correction is 76.6.
This concrete example shows why a familiar benchmark should not be treated as unquestionable ground truth.
A corrected-data comparison would be a separately identified analysis, with potentially different results.

\section{Load the table and understand its scales}
The small public CSV is downloaded with a 20-second timeout. The first column becomes the DataFrame index;
only the four named numeric columns are converted to the feature matrix.
Checks require shape $(50,4)$, unique identifiers, finite numbers, and no missing values.
These checks detect common import failures, but do not establish the validity of the historical measurements.

\begin{center}\begin{tabular}{lrrr}\toprule
Feature&Mean&Sample SD&Sample variance\\\midrule
Murder&7.788&4.3555&18.9705\\
Assault&170.760&83.3377&6945.1657\\
UrbanPop&65.540&14.4748&209.5188\\
Rape&21.232&9.3664&87.7292\\\bottomrule\end{tabular}\end{center}
Assault's raw variance is much larger than the others. A one-unit change in an arrest rate and a one-percentage-point change
in urban population have different meanings; adding their squared differences is a modeling decision.
The large Assault mean is not the reason raw PCA emphasizes it: PCA removes all means before measuring variation.

\code{read_csv(index_col=0)} removes identifiers from the numeric matrix at import.
\code{std(ddof=1)} and \code{var(ddof=1)} use a divisor of $n-1$; \code{ddof} is the number subtracted from the sample size.
Columns should be selected by name before converting to an array so their ordering remains explicit throughout interpretation.

\section{Relate standardized PCA to the correlation matrix}
Let $s_j$ be the sample standard deviation and $Z_{ij}=(X_{ij}-\bar x_j)/s_j$.
Each standardized column has sample variance 1, so $Z^TZ/(n-1)$ is the sample correlation matrix and its trace is 4.
Eigenvalues sum to that trace. The notebook checks this identity independently of the displayed variance ratios.

\begin{center}\begin{tabular}{rrrr}\toprule
PC&Raw PVE&Standardized PVE&Standardized CPVE\\\midrule
1&0.965534&0.620060&0.620060\\
2&0.027817&0.247441&0.867502\\
3&0.005800&0.089141&0.956642\\
4&0.000849&0.043358&1.000000\\\bottomrule\end{tabular}\end{center}
PVE is proportion of variance explained; CPVE is its cumulative sum.
The raw first direction has coefficients approximately $(0.0417,0.9952,0.0463,0.0752)$,
so its very high PVE mainly summarizes Assault variation in the chosen raw units.
Standardization distributes numerical emphasis more evenly, producing a less concentrated variance spectrum.
Neither percentage is a classification accuracy.

The first two standardized direction coefficients are:
\begin{center}\begin{tabular}{lrr}\toprule
Feature&PC1&PC2\\\midrule
Murder&0.535899&$-0.418181$\\
Assault&0.583184&$-0.187986$\\
UrbanPop&0.278191&0.872806\\
Rape&0.543432&0.167319\\\bottomrule\end{tabular}\end{center}
The sign convention makes the largest-magnitude coefficient in each direction positive.
PC1 combines positively associated variation across the four variables, with stronger direction coefficients on the arrest-rate features.
PC2 mainly contrasts higher urban-population percentage with lower values along the Murder and Assault contributions, conditional on this coordinate description.
It is a linear contrast, not evidence that urbanization causes any change in arrests.

\subsection*{Direction coefficients are not automatically correlations}
The notebook deliberately labels its table as direction coefficients. For standardized feature column $Z_j$ and score $t_\ell=Zv_\ell$,
\[
\operatorname{Cov}(Z_j,t_\ell)=(Sv_\ell)_j=\lambda_\ell v_{j\ell}.
\]
Because $\operatorname{Var}(Z_j)=1$ and $\operatorname{Var}(t_\ell)=\lambda_\ell$,
\[\operatorname{Corr}(Z_j,t_\ell)=\sqrt{\lambda_\ell}\,v_{j\ell}.\]
Thus a coefficient such as 0.5359 is not itself a correlation. Different sources use the word ``loading'' for different quantities;
read their definition before comparing tables.

\section{Interpret the score plot and reconstruction}
Scores give state coordinates. Multiplying a direction and its scores by the same sign leaves every reconstruction unchanged.
Mixing signed scores with unsigned directions, however, would be an implementation error.
The notebook uses the original fitted directions for reconstruction and the sign-adjusted coordinates only for interpretation and plotting.

\fig{04_C_usarrests_1.png}{The two-PC score plot describes 86.75\% of standardized variation. Nearby plotted states may still differ in the discarded coordinates.}
The standardized reconstruction is
\[\widehat Z=\bar z+(Z-\bar z)V_2V_2^T,\]
where the PCA fitting mean $\bar z$ is numerically close to zero. Original units are restored by
$\widehat X_j=\bar x_j+s_j\widehat Z_j$.
The sum of squared standardized residuals is $49(\lambda_3+\lambda_4)$, about 25.97.
Per-feature root mean squared errors (RMSEs), after restoring original units, are:
\begin{center}\begin{tabular}{lr}\toprule
Feature&Two-PC reconstruction RMSE\\\midrule
Murder&1.4598\\Assault&28.7552\\UrbanPop&3.3317\\Rape&4.5408\\\bottomrule\end{tabular}\end{center}
These values quantify approximation of the observed table. Assault's numerically largest RMSE is partly a consequence of its units.
PCA after scaling minimizes weighted squared errors $\sum_{i,j}(X_{ij}-\widehat X_{ij})^2/s_j^2$,
not the unweighted sum in original units. The residual summaries therefore belong to a clearly specified optimization geometry.

\section{Select K before interpreting the named groups}
For each K from 2 to 6, run K-means with 50 starts and keep the lowest inertia. Select the largest mean silhouette in the full
four-dimensional standardized space, breaking ties toward smaller K.
\begin{center}\begin{tabular}{rrr}\toprule
K&Inertia&Silhouette\\\midrule
2&102.8624&0.408489\\3&78.4840&0.308136\\4&56.4032&0.339689\\
5&48.9442&0.303078\\6&42.9639&0.288293\\\bottomrule\end{tabular}\end{center}
The chosen K is 2, under this criterion and candidate set. It is not an estimate of a universally true number of types of states.
The fitted profiles in original units are:
\begin{center}\begin{tabular}{rrrrrr}\toprule
ID&Size&Murder&Assault&UrbanPop&Rape\\\midrule
0&30&4.870&114.433&63.633&15.943\\
1&20&12.165&255.250&68.400&29.165\\\bottomrule\end{tabular}\end{center}
These describe higher versus lower values of the measured arrest variables, with a smaller urban-population difference.
A cluster ID is not an ordinal score. Renaming group 1 as group 0 changes no mathematical result.
The groups provide no individual-level inference and no statement about present-day states.

Because scaling is an affine transformation, inverse-transforming a standardized centroid gives the corresponding original-unit group mean.
That identity does not imply that running K-means directly on unscaled X would produce these groups.
The latter procedure would minimize a different objective and emphasize different features.

\section{Compare methods in a common evaluation space}
At the fixed selected K, Ward and average hierarchies use all four standardized features.
Another K-means fit uses only the first two scores. All silhouettes in this comparison are evaluated using full Z,
so differences in their reported values do not arise merely from changing the scoring geometry.
\begin{center}\begin{tabular}{lrr}\toprule
Method&ARI versus full K-means&Full-space silhouette\\\midrule
Full K-means&1.000000&0.408489\\Ward&0.919937&0.404794\\
Average&1.000000&0.408489\\Two-PC K-means&1.000000&0.408489\\\bottomrule\end{tabular}\end{center}
The two-PC partition happens to agree exactly here. It is important to report that observation instead of claiming dimension reduction
must change groups. Truncation does change some distances, but this particular optimum remains in the same assignment region.
Ward's greedy hierarchy differs slightly, illustrating that related squared-error ideas need not yield identical final partitions.

For a full orthogonal basis V, $(ZV)(ZV)^T=ZZ^T$, so all centered row inner products and Euclidean distances are preserved.
The notebook checks this Gram-matrix identity. After truncation, the omitted part is $ZV_\perp V_\perp^TZ^T$;
there is no general distance-preservation guarantee. A two-PC plot cannot alone settle the quality of full-space groups.

\fig{04_C_usarrests_2.png}{Ward hierarchy in standardized Euclidean geometry. Heights equal the square root of twice the added within-group squared error. Horizontal leaf ordering is a drawing choice.}

\section{What the stability experiment establishes}
Ten 40-of-50 subsamples each recompute their own mean, sample standard deviation, and K-means fit at fixed K=2.
Pairwise overlaps range from 30 to 35 states, with median 32. All 45 shared-row ARIs are 1 in this run.
Thus the partition is stable under these particular deletions despite the slight method disagreement above.
Sampling sensitivity and choice-of-method sensitivity are different questions.

Keeping the original scaler for all subsets would assess a fixed-metric clustering procedure; refitting it assesses the combined
preprocessing-and-clustering procedure used here. Neither choice is universally mandatory, but the experiment must identify which it uses.
K is held fixed, so the result does not assess how often silhouette would reselect the same K.
Pairwise scores share fits and should not be treated as 45 independent samples for a standard error.
Finally, the historical states are not naturally independent interchangeable draws from a simple population;
this is a sensitivity analysis of the table, not a forecast confidence statement.

\section*{Documentation and practical tips}
Read the \href{https://scikit-learn.org/1.8/modules/generated/sklearn.decomposition.PCA.html}{PCA API}
for direction-row orientation and sample covariance conventions, and the
\href{https://scikit-learn.org/1.8/modules/generated/sklearn.cluster.KMeans.html}{KMeans API}
for restart and learned-center attributes. Its \code{transform} returns distances to centers, whereas PCA's \code{transform}
returns principal-component coordinates. Identically named methods can have different mathematical outputs.

The notebook manually uses sample scaling. \code{StandardScaler} uses population standard deviations, smaller by
$\sqrt{49/50}$ here, so its standardized coordinates are larger by $\sqrt{50/49}$.
Directions, PVE, silhouettes, and exact K-means minimizing partitions are unchanged by this common positive factor;
eigenvalues and inertia multiply by $50/49$. Solver stopping tolerances can still produce tiny numerical differences.

Label every plot's coordinate space, retain state identifiers outside numerical features, and keep known data-quality limitations in the interpretation.
Similarity of aggregate records is not evidence of a causal relation between the measured variables.
\end{document}
